% Terminal apparatus of the concise study: the scope note and the References
% for the sources this companion actually uses, at the loci it cites. The
% revision timestamp and the rights colophon follow from the shared
% generation record.

\runninghead{Scope and References}
\section{Appendix: Scope and Qualifications}

{\footnotesize
\textbf{Edition, Mass and occurrence.} The Mass treated is \latin{Dominica
sexta quae superfuit post Epiphaniam}, II classis, of the \work{Missale
Romanum}, editio typica, Typis Polyglottis Vaticanis 1962, printed
pp.~421--422, nos.~1759--1768, as said on Sunday 15~November 2026, the
twenty-fifth Sunday after Pentecost, in the universal 1962 calendar. Its
orations and lessons are those of the Sixth Sunday after the Epiphany and its
chants those of the Twenty-third Sunday after Pentecost, under RGMR~298. Its
place in 2026, the omission of the calendar feast of St Albert the Great, the
colour, the Gloria, the Credo and the Trinity Preface rest on the Missal's
general rubrics and on its printed direction after the Secret, as
\texttt{research/context.md} records. The Ordinary and the Office are not
treated.

\textbf{Text and English.} The Latin of the ten texts was read on page images
of the CMAA facsimile of the 1962 typical edition and collated with the
Ratisbon printing of 1862 (Pustet); the collation is in
\texttt{propers/verified.md}, and the expansive study prints every text whole.
The English of the chants and orations is that of the Lasance \work{New Roman
Missal} (Benziger, 1945), public domain in the United States by non-renewal,
which prints this Mass as its \work{Supplementary Mass No.~1}; the lessons are
quoted in the Douay--Rheims in Challoner's revision. None of this English is
an approved liturgical translation of the 1962 Missal. Augustine on the Psalms
and Chrysostom on Matthew and on First Thessalonians are quoted in the English
of the Nicene and Post-Nicene Fathers; the expositions of Psalms~43, 77 and 84
in that series are abridged, so that the quotations are of the condensed
English. The other Latin of the Fathers is given with the English rendering
made for the expansive study; no English is composed for a liturgical text.

\textbf{Reception.} The three interpretations are developed in full in the
expansive study; their loci, editions and bounds are recorded in
\texttt{research/scope.md} and \texttt{research/interpretations.md}. Every
connection drawn between two texts of the Mass belongs to the interpretations
unless an author is named as making it; no Father cited here comments on this
Mass as a whole. Ambrose, Chrysologus and Bede are cited on Luke's form of the
two parables, Bede also on Mark's verses of the Communion's saying, and
Chrysostom on Matthew's form of it (Mt 21:21--22). The Latin of Jerome and of
Hilary on Matthew was read on page images; Ambrose and Bede in transcriptions
of Migne not collated with page images;
Hilary on the Psalms, Chrysologus, Aquinas, Rupert, Durandus and Sicard in
uncorrected optical or web text; Chrysostom on the Psalms and Theodoret in
uncorrected Greek optical text, so that they are given in paraphrase. Rupert,
Durandus and Sicard are cited for the chants alone. Theodoret, Rupert,
Durandus and Sicard are ecclesiastical writers and are cited in support only.
No commentary on the Collect, Secret or Postcommunion survives among the
commentators read.

\textbf{Chronology.} Every biblical date on page~2 is generated from the
repository's Scripture chronology corpus; the expansive study's Scriptural
Date and Location appendix gives each passage's sources and evidence.

\textbf{Rights and use.} The Latin is published on 17 U.S.C. 103(b) together
with its public-domain witness, the Pustet printing of 1862. The Douay--Rheims
and the Lasance English are quoted from registered editions. The other
quotations are short and come from editions published before 1931, named
below, or from medieval and ancient Latin; the Corpus Thomisticum and
Wikisource texts are used only for Latin that is itself medieval or ancient.
Scripture, liturgical texts and quoted translations are not relicensed by the
repository's licence. This is a study aid to the 1962 Missal, not an official
liturgical book, translation or critical edition.
\par}

\section{References}

{\footnotesize
\begin{itemize}[leftmargin=1.2em,itemsep=.15em]
\item \work{Missale Romanum ex decreto Sacrosancti Concilii Tridentini
restitutum, Summorum Pontificum cura recognitum}, editio typica (Typis
Polyglottis Vaticanis, 1962), CMAA facsimile: \work{Rubricae generales} and
\work{Rubricae generales Missalis romani}; Proprium de Tempore,
pp.~421--422, nos.~1759--1768.
\item \work{Missale Romanum} (Ratisbon: Pustet, 1862), pp.~44--45 and
357--359.
\item Lasance, F.~X., and Walsh, F.~A., \work{The New Roman Missal in Latin and
English}, revised edition (New York: Benziger Brothers, 1945), pp.~740--743.
\item \work{The Holy Bible Translated from the Latin Vulgate} (Douay--Rheims,
Challoner revision), Project Gutenberg text: Jer 29:5--14; Ps 43:4--10, 22;
84:2--5; 129:1--2; Mt 13:31--36; 16:6; Mk 11:23--25; Acts 17:1--10;
1~Cor 5:7; 1~Thess 1:2--10.
\item Ambrose, \work{Expositio evangelii secundum Lucam} VII.176--194, on the
mustard seed and the leaven, Wikisource transcription of PL~15.
\item Augustine, \work{Enarrationes in Psalmos} 43, 77 and 84, NPNF, first
series, vol.~8; \work{Quaestiones Evangeliorum} I.11--12, Migne Latin text
(Wikisource).
\item Bede, \work{In Lucae evangelium expositio} IV, on Lk 13:19--21, PL~92,
Wikisource transcription; \work{In Marci evangelium expositio} III, on
11:22--25, Migne Latin text (Wikisource).
\item Chrysologus, Peter, \work{Sermones} 98 and 99, PL~52 (Paris, 1894).
\item Chrysostom, John, \work{Homilies on Matthew} 46, 47 and 67, NPNF, first
series, vol.~10; \work{Homilies on First Thessalonians} 1 and 2, NPNF, first
series, vol.~13; \work{Expositio in Psalmos}, on Pss 43 and 129, PG~55.
\item Hilary of Poitiers, \work{In Evangelium Matthaei commentarius} 13.4--6,
PL~9 (Paris, 1844), cols.~994--995; \work{Tractatus super Psalmos}, on Ps 129,
PL~9, col.~718.
\item Jerome, \work{Commentarii in Matthaeum} II, on 13:31--35, PL~26 (Paris,
1845), cols.~90--93; \work{In Hieremiam} V, on 29:10--14, PL~24 (Paris,
1865), cols.~893--894.
\item Theodoret of Cyrus, \work{Interpretatio epistolae I ad Thessalonicenses},
on 1:3, PG~82.
\item Thomas Aquinas, \work{Super Evangelium S.~Matthaei lectura}, cap.~13, in
\work{Opera omnia}, vol.~3 (Venice: Bettinelli, 1745), pp.~186--189;
\work{Super I epistolam ad Thessalonicenses}, cap.~1, lect.~1, Corpus
Thomisticum text; \work{Super Psalmos}, on Ps 43, n.~3, Corpus Thomisticum
text.
\item Durandus, William, \work{Rationale divinorum officiorum} VI.140, in the
Lyon printing of 1612, tomus secundus, p.~412.
\item Rupert of Deutz, \work{De divinis officiis} XII.23, PL~170,
cols.~329--330.
\item Sicard of Cremona, \work{Mitrale} VIII.23, PL~213.
\end{itemize}
\par}
